\section{Perturbation Rendering of the Mandelbrot Set}
\label{sec:perturbation-concept}

Rendering the Mandelbrot set (\cref{sec:mandelbrot-intro}) at extreme zoom levels presents a
numerical challenge. As the image is magnified, the parameters \(c\) associated
with individual pixels differ by increasingly small amounts, while the number
of iterations required to determine escape behavior typically grows. Accurately
tracking these orbits therefore requires both high numerical precision and a
large number of iterations, making naive per-pixel evaluation
expensive.

Perturbation provides a mathematical and algorithmic framework
\cite{Martin2013Perturbation,HeilandAllen2013Perturbation,Wikibooks_Perturbation}
that addresses
this challenge by exploiting the strong coherence between nearby orbits. Rather
than evaluating each pixel independently at full precision, perturbation
separates the computation into a shared high-precision component and many
pixel-local low-precision components. This section introduces the perturbation
approach and explains how it enables efficient deep-zoom Mandelbrot rendering.

\subsection{Baseline iteration and its limitations}

The Mandelbrot set is defined by the iteration
\begin{equation}
z_{n+1} = z_n^2 + c, \qquad z_0 = 0,
\label{eq:mandelbrot-iter}
\end{equation}
where \(c \in \mathbb{C}\) is the parameter associated with a pixel and the orbit
\(\{z_n\}\) is iterated until either escape is detected or a maximum iteration
count is reached.

At modest zoom levels, this iteration can be evaluated accurately using standard
floating-point arithmetic. At deep zoom levels, however, the differences between
nearby parameters \(c\) may be many orders of magnitude smaller than the values
of \(z_n\) encountered during iteration. In such regimes, rounding error and
loss of significance make fixed-precision arithmetic unreliable. While arbitrary
precision arithmetic can restore accuracy, applying it independently to every
pixel scales poorly.

\subsection{Coherence of nearby orbits}

\begin{figure}[!htbp]
\centering
\begin{tikzpicture}[
    >=Stealth,
    scale=0.85,
    every node/.style={font=\small},
]
  % Image tile with pixel grid
  \draw[thick] (0,0) rectangle (5,5);
  \node[font=\scriptsize, above] at (2.5, 5) {Image tile (parameter plane)};

  % Faint pixel grid
  \foreach \x in {0.5,1,...,4.5} {
    \draw[gray!15] (\x, 0) -- (\x, 5);
    \draw[gray!15] (0, \x) -- (5, \x);
  }

  % Concentric distance rings labeled with R fractions
  \draw[green!60!black, dashed] (2.5, 2.5) circle (0.6);
  \node[green!60!black, font=\tiny, anchor=south west] at (2.9, 3.0) {$\tfrac{1}{4}R$};
  \draw[orange!70!black, dashed] (2.5, 2.5) circle (1.4);
  \node[orange!70!black, font=\tiny, anchor=south west] at (3.4, 3.7) {$\tfrac{1}{2}R$};
  \draw[red!60!black, dashed] (2.5, 2.5) circle (2.2);
  \node[red!60!black, font=\tiny, anchor=south west] at (4.0, 4.4) {$R$};

  % Reference point at center
  \fill[blue!70!black] (2.5, 2.5) circle (4pt);
  \node[blue!70!black, font=\scriptsize, below=4pt] at (2.5, 2.5) {$c_\star$};

  % Close pixels (small delta) — green tint
  \fill[green!60!black] (2.9, 2.8) circle (2pt);
  \fill[green!60!black] (2.1, 2.9) circle (2pt);
  \fill[green!60!black] (2.7, 2.1) circle (2pt);

  % Medium pixels
  \fill[orange!80!black] (1.5, 3.5) circle (2pt);
  \fill[orange!80!black] (3.6, 1.5) circle (2pt);
  \fill[orange!80!black] (3.5, 3.7) circle (2pt);

  % Far pixels (large delta) — red tint
  \fill[red!70!black] (0.5, 4.5) circle (2pt);
  \fill[red!70!black] (4.5, 0.5) circle (2pt);
  \fill[red!70!black] (4.5, 4.2) circle (2pt);

  % Right side: delta magnitude vs distance
  \begin{scope}[xshift=7cm, yshift=0.3cm]
    \draw[->, gray] (0, 0) -- (0, 4.2)
      node[above, font=\scriptsize] {$|\Delta z_n|$};
    \draw[->, gray] (0, 0) -- (4.5, 0)
      node[right, font=\scriptsize] {$|\Delta c|$};

    % R tick marks on x-axis
    \draw[green!60!black] (1.0, -0.1) -- (1.0, 0.1);
    \node[green!60!black, font=\tiny, below] at (1.0, -0.1) {$\tfrac{1}{4}R$};
    \draw[orange!70!black] (2.0, -0.1) -- (2.0, 0.1);
    \node[orange!70!black, font=\tiny, below] at (2.0, -0.1) {$\tfrac{1}{2}R$};
    \draw[red!60!black] (3.8, -0.1) -- (3.8, 0.1);
    \node[red!60!black, font=\tiny, below] at (3.8, -0.1) {$R$};

    % Roughly linear growth curve
    \draw[thick, blue!60!black]
      plot[smooth, domain=0:3.8] (\x, {0.8*\x});

    % Example pixels on curve
    \fill[green!60!black] (1.0, 0.8) circle (2.5pt);
    \node[green!60!black, font=\tiny, above left] at (1.0, 0.8)
      {$\tfrac{1}{4}R$};

    \fill[orange!80!black] (2.0, 1.6) circle (2.5pt);
    \node[orange!80!black, font=\tiny, above left] at (2.0, 1.6)
      {$\tfrac{1}{2}R$};

    \fill[red!70!black] (3.8, 3.04) circle (2.5pt);
    \node[red!70!black, font=\tiny, above left] at (3.8, 3.04)
      {$R$};

    % Precision boundary
    \draw[dashed, purple!60!black, thick] (-0.2, 2.0) -- (4.3, 2.0);
    \node[purple!60!black, font=\tiny, right, align=left] at (4.3, 2.0)
      {precision\\limit};
  \end{scope}

  % Bottom labels
  \node[blue!70!black, font=\scriptsize, align=center] at (2.5, -0.5)
    {$R = \code{MaxRadius}$ (view half-height)};
  \node[font=\scriptsize, align=center] at (8.5, -0.5)
    {Per-pixel $|\Delta z_n|$ grows with $|\Delta c|$};
\end{tikzpicture}
\caption[Perturbation decomposition: spatial delta growth]{%
  Perturbation decomposition.  The reference orbit at $c_\star$
  (blue, center) is computed once at full precision.  Each pixel's
  perturbation delta $\Delta z_n$ grows with both the spatial distance
  $|\Delta c|$ from the reference and the number of iterations elapsed.
  Dashed circles mark fractions of $R = \code{MaxRadius}$, the view
  half-height in the parameter plane.  When $|\Delta z_n|$ exceeds the
  working precision (purple line), the pixel must be \emph{rebased}
  to recover accuracy---a common occurrence even for pixels well
  inside the view.}
\label{fig:perturbation-delta}
\end{figure}

A key observation underlying perturbation is that nearby parameters generate
orbits that remain close for many iterations, as illustrated in
\cref{fig:perturbation-delta}. In a typical image tile, all pixel
parameters \(c\) lie within a small neighborhood of a central value. Their
orbits therefore share a common large-scale structure, differing only by small
corrections that grow gradually over time.

Perturbation makes this structure explicit by selecting a single
\emph{reference parameter} and expressing all nearby orbits relative to it. This
transforms the problem from one of many independent high-precision computations
into one of a single high-precision computation plus many low-precision updates.

\subsection{Reference orbit}

Let \(c_\star\) denote a reference parameter, typically chosen as the center of
the image region. Its orbit is computed using sufficiently high precision to
serve as a reliable baseline:
\begin{equation}
z_{n+1}^\star = (z_n^\star)^2 + c_\star, \qquad z_0^\star = 0.
\label{eq:ref-orbit-recurrence}
\end{equation}
The sequence \(\{z_n^\star\}\) is stored and reused for all pixels within the
region. Because this orbit underpins the evaluation of many pixels, its
accuracy is critical.

\subsection{Perturbation as a mixed-precision decomposition}
\label{subsec:perturbation-mixed-precision}

For a nearby pixel parameter \(c = c_\star + \Delta c\), the orbit can be written
as a small deviation from the reference orbit:
\begin{equation}
z_n = z_n^\star + \Delta z_n.
\label{eq:perturb-decomposition}
\end{equation}
Substituting this expression into the Mandelbrot recurrence yields an exact
update rule for the deviation:
\begin{equation}
\Delta z_{n+1}
= 2 z_n^\star \Delta z_n + (\Delta z_n)^2 + \Delta c.
\label{eq:perturb-exact}
\end{equation}
No approximation has been introduced: as long as the reference orbit is computed
accurately, this formulation preserves the exact Mandelbrot dynamics.

This decomposition has an important numerical consequence. The reference orbit
\(z_n^\star\) typically grows in magnitude and requires arbitrary precision to be
represented accurately at deep zoom. In contrast, the deviation \(\Delta z_n\)
remains small for many iterations and can usually be evolved using much lower
precision arithmetic. Perturbation therefore reorganizes the computation into a
mixed-precision pipeline: expensive high-precision arithmetic is concentrated in
a single shared orbit, while the majority of per-pixel work is performed using
lower precision with significantly higher throughput.

\subsection{Real-valued formulation for implementation}

For practical implementation, particularly on GPUs, it is convenient to express
the perturbation update in real arithmetic. Writing
\begin{align*}
z_n^\star &= x_n^\star + i y_n^\star, \\
\Delta z_n &= \Delta x_n + i \Delta y_n, \\
\Delta c &= \Delta c_x + i \Delta c_y,
\end{align*}
and expanding \cref{eq:perturb-exact} yields
\begin{align}
\Delta x_{n+1}
&= 2(x_n^\star \Delta x_n - y_n^\star \Delta y_n)
   + (\Delta x_n^2 - \Delta y_n^2)
   + \Delta c_x, \\
\Delta y_{n+1}
&= 2(x_n^\star \Delta y_n + y_n^\star \Delta x_n)
   + 2 \Delta x_n \Delta y_n
   + \Delta c_y.
\end{align}
These equations map directly onto scalar arithmetic and can be evaluated
efficiently in parallel for many pixels.

\subsection{Reconstruction and escape testing}

Although perturbation evolves only the deviation \(\Delta z_n\), escape-time
rendering requires testing the magnitude of the full orbit value. At each
iteration, the current iterate is reconstructed as
\[
z_n = z_n^\star + \Delta z_n,
\]
and the standard Mandelbrot bailout condition is applied:
\[
|z_n|^2 \ge 4.
\]
The iteration at which this condition is first satisfied determines the pixel's
escape time.

\subsection{Numerical stability and rebasing}

Perturbation remains efficient only while the deviation \(\Delta z_n\) remains
small relative to the reference orbit. As the deviation grows, numerical
cancellation and loss of significance may occur. Practical implementations
therefore employ \emph{rebasing} \cite{Zhuoran2023ReferenceCompression,HeilandAllen2022DeepZoom},
in which the current deviation
is folded into the reference state and perturbation continues relative to this
new baseline. This process preserves accuracy while maintaining numerical
stability.

\subsection{Relation to other acceleration techniques}

Perturbation preserves the exact dynamics of the Mandelbrot map and forms the
foundation of many advanced deep-zoom rendering techniques. It can be used by
itself (\cref{sec:perturb-only}) or combined with additional approximations, such as bilinear approximation
(\cref{sec:bilinear-approx}) or linear approximation (\cref{sec:la-v2-perturb}),
to skip multiple iterations at once. In all cases, perturbation provides
the conceptual framework that makes efficient deep-zoom Mandelbrot rendering
possible.

\paragraph{Terminology note: HDR scalar representation.}
In this work, the term \emph{HDR} refers to a \emph{high-dynamic-range} numeric
representation that pairs a fixed-precision mantissa (float or double) with an
explicitly managed integer exponent (see below for the full description).  When the GPU reference-orbit kernel or
perturbation pipeline needs lightweight control-flow decisions---periodicity
detection, escape testing, magnitude comparisons---it converts the full
arbitrary-precision state into an HDR value via a lossy reduction step
(\code{ToHDRFloat}).  The resulting HDR value preserves a wide exponent range,
sufficient to compare magnitudes spanning many orders of magnitude, but carries
only the fixed mantissa precision of the underlying float type rather than the
full arbitrary-precision mantissa.  The use of the term HDR here is purely
numerical and is unrelated to high-dynamic-range imaging.
